\section{Banco de questões — Direito Administrativo}

\subsection*{N1 — Fundamento competitivo}

\begin{questao}{\nivel{N1} 04.01 — Cebraspe/Certo ou Errado}
Em sentido subjetivo, Administração Pública designa os órgãos, agentes e entidades que exercem função administrativa; em sentido objetivo, refere-se à própria atividade administrativa.
\end{questao}

\begin{questao}{\nivel{N1} 04.02 — Cebraspe/Certo ou Errado}
A criação de novos órgãos dentro de uma mesma pessoa jurídica caracteriza descentralização administrativa, pois distribui competências entre centros distintos.
\end{questao}

\begin{questao}{\nivel{N1} 04.03 — Cebraspe/Certo ou Errado}
Órgãos públicos, embora possuam competências próprias, não dispõem de personalidade jurídica autônoma em relação à pessoa jurídica que integram.
\end{questao}

\begin{questao}{\nivel{N1} 04.04 — Cebraspe/Certo ou Errado}
Todo ato administrativo possui, necessariamente, imperatividade e autoexecutoriedade.
\end{questao}

\begin{questao}{\nivel{N1} 04.05 — Cebraspe/Certo ou Errado}
A Administração pode anular ato ilegal e revogar ato válido por razões de conveniência e oportunidade, observados os limites jurídicos aplicáveis.
\end{questao}

\begin{questao}{\nivel{N1} 04.06 — Cebraspe/Certo ou Errado}
O Poder Judiciário, no exercício típico do controle jurisdicional, pode revogar ato administrativo válido sempre que considerar inconveniente a escolha feita pela Administração.
\end{questao}

\begin{questao}{\nivel{N1} 04.07 — Cebraspe/Certo ou Errado}
Efetividade e estabilidade são conceitos equivalentes: todo servidor ocupante de cargo efetivo é estável desde a posse.
\end{questao}

\begin{questao}{\nivel{N1} 04.08 — Cebraspe/Certo ou Errado}
A supervisão finalística exercida sobre entidade da Administração indireta não se confunde com hierarquia existente entre órgãos da mesma pessoa jurídica.
\end{questao}

\begin{questao}{\nivel{N1} 04.09 — Cebraspe/Certo ou Errado}
Na inexigibilidade de licitação, a característica central é a inviabilidade de competição; na dispensa, a lei admite contratação direta em hipóteses previstas mesmo quando competição poderia existir.
\end{questao}

\begin{questao}{\nivel{N1} 04.10 — Cebraspe/Certo ou Errado}
Após as alterações promovidas pela Lei nº 14.230/2021, a mera culpa do agente continua suficiente, em qualquer hipótese, para caracterizar ato de improbidade administrativa.
\end{questao}

\subsection*{N2 — Aplicação}

\begin{questao}{\nivel{N2} 04.11 — FGV/organização administrativa}
Um ministério cria internamente uma secretaria especializada, sem instituição de nova pessoa jurídica. O fenômeno é:
\begin{enumerate}[label=\Alph*)]
\item descentralização por serviços;
\item desconcentração administrativa;
\item privatização;
\item delegação contratual de serviço público;
\item criação de autarquia.
\end{enumerate}
\end{questao}

\begin{questao}{\nivel{N2} 04.12 — FCC/administração indireta}
Assinale a alternativa correta.
\begin{enumerate}[label=\Alph*)]
\item Autarquia é órgão sem personalidade jurídica.
\item Empresa pública possui necessariamente capital misto público-privado.
\item Sociedade de economia mista integra a Administração indireta e possui personalidade jurídica própria.
\item Fundação pública nunca integra a Administração indireta.
\item Toda entidade indireta está subordinada hierarquicamente ao ministério supervisor.
\end{enumerate}
\end{questao}

\begin{questao}{\nivel{N2} 04.13 — Cebraspe/Certo ou Errado}
A discricionariedade administrativa existe apenas dentro dos espaços conferidos pelo ordenamento e continua sujeita a controle de finalidade, motivos, proporcionalidade e juridicidade.
\end{questao}

\begin{questao}{\nivel{N2} 04.14 — FGV/atos administrativos}
A autoridade identifica vício de competência em ato praticado por agente incompetente, mas a competência não era exclusiva e não houve prejuízo ao interesse público nem a terceiros. Em tese, a providência juridicamente possível é:
\begin{enumerate}[label=\Alph*)]
\item revogação obrigatória;
\item convalidação, se preenchidos os requisitos legais;
\item prescrição da competência;
\item conversão automática em lei;
\item declaração de inexistência do ato em qualquer caso.
\end{enumerate}
\end{questao}

\begin{questao}{\nivel{N2} 04.15 — FCC/anulação e revogação}
Um ato válido deixou de ser conveniente para a Administração, sem gerar direito adquirido à sua manutenção. A medida típica é:
\begin{enumerate}[label=\Alph*)]
\item anulação por ilegalidade;
\item revogação por mérito administrativo;
\item convalidação;
\item cassação judicial obrigatória;
\item declaração de nulidade absoluta pelo particular.
\end{enumerate}
\end{questao}

\begin{questao}{\nivel{N2} 04.16 — Cebraspe/Certo ou Errado}
A presunção de legitimidade dos atos administrativos é relativa e admite prova em contrário.
\end{questao}

\begin{questao}{\nivel{N2} 04.17 — FGV/poderes administrativos}
Um agente competente para aplicar determinada sanção utiliza essa atribuição para retaliar servidor por divergência pessoal. O vício predominante caracteriza:
\begin{enumerate}[label=\Alph*)]
\item excesso de poder;
\item desvio de finalidade;
\item delegação válida;
\item discricionariedade legítima;
\item convalidação tácita.
\end{enumerate}
\end{questao}

\begin{questao}{\nivel{N2} 04.18 — FCC/poder hierárquico}
O poder hierárquico permite, observados os limites legais:
\begin{enumerate}[label=\Alph*)]
\item criar hierarquia entre autarquia e ministério supervisor;
\item distribuir funções, fiscalizar subordinados e rever atos internos;
\item aplicar sanção penal a qualquer cidadão;
\item editar lei formal;
\item afastar controle judicial.
\end{enumerate}
\end{questao}

\begin{questao}{\nivel{N2} 04.19 — Cebraspe/Certo ou Errado}
Excesso de poder relaciona-se ao ultrapassamento da competência; desvio de finalidade, ao uso da competência para finalidade diversa da prevista juridicamente.
\end{questao}

\begin{questao}{\nivel{N2} 04.20 — FGV/responsabilidade civil}
Veículo oficial, conduzido por servidor no exercício da função, causa dano a terceiro por atuação do agente. À luz do art. 37, §6º, da Constituição, a responsabilidade da pessoa jurídica é, em regra:
\begin{enumerate}[label=\Alph*)]
\item subjetiva, exigindo prova de dolo do Estado;
\item objetiva, sem prejuízo do direito de regresso nos casos cabíveis;
\item inexistente, pois o agente responde sozinho;
\item penal;
\item exclusivamente contratual.
\end{enumerate}
\end{questao}

\begin{questao}{\nivel{N2} 04.21 — Cebraspe/Certo ou Errado}
A responsabilidade objetiva do Estado dispensa a demonstração de nexo causal entre a atuação estatal e o dano.
\end{questao}

\begin{questao}{\nivel{N2} 04.22 — FCC/agentes públicos}
A estabilidade constitucional do servidor ocupante de cargo efetivo:
\begin{enumerate}[label=\Alph*)]
\item decorre imediatamente da nomeação;
\item confunde-se com efetividade;
\item depende do atendimento aos requisitos constitucionais próprios e não nasce automaticamente com a posse;
\item alcança todo ocupante de cargo em comissão;
\item é idêntica à vitaliciedade.
\end{enumerate}
\end{questao}

\begin{questao}{\nivel{N2} 04.23 — FGV/Lei 9.784}
A autoridade pretende avocar competência de órgão inferior por tempo indeterminado e sem justificativa. A medida é:
\begin{enumerate}[label=\Alph*)]
\item compatível com a lei, pois avocação é livre;
\item incompatível, pois a avocação é excepcional, temporária e deve ser justificada nos termos legais;
\item obrigatória quando houver hierarquia;
\item equivalente à delegação;
\item possível apenas por decisão judicial.
\end{enumerate}
\end{questao}

\begin{questao}{\nivel{N2} 04.24 — Cebraspe/Certo ou Errado}
Na Lei nº 9.784/1999, determinadas matérias, como decisão de recursos administrativos e edição de atos de caráter normativo, não podem ser objeto de delegação.
\end{questao}

\begin{questao}{\nivel{N2} 04.25 — FCC/licitações}
Na Lei nº 14.133/2021, contratação direta por inviabilidade de competição enquadra-se como:
\begin{enumerate}[label=\Alph*)]
\item dispensa;
\item inexigibilidade;
\item pregão obrigatório;
\item concorrência simplificada;
\item diálogo competitivo necessariamente.
\end{enumerate}
\end{questao}

\subsection*{N3 — Integração de concurso}

\begin{questao}{\nivel{N3} 04.26 — FGV/assertivas sobre atos}
Considere:

I. A anulação decorre de ilegalidade.

II. A revogação pressupõe ato válido e juízo de conveniência/oportunidade.

III. O Judiciário pode revogar ato administrativo válido no controle típico de legalidade.

Assinale a alternativa correta.
\begin{enumerate}[label=\Alph*)]
\item Apenas I.
\item Apenas II.
\item Apenas I e II.
\item Apenas II e III.
\item I, II e III.
\end{enumerate}
\end{questao}

\begin{questao}{\nivel{N3} 04.27 — Cebraspe/Certo ou Errado}
A decadência administrativa prevista na Lei nº 9.784/1999 pode limitar o poder de anular atos favoráveis ao destinatário, ressalvada a hipótese de má-fé e demais condições legais.
\end{questao}

\begin{questao}{\nivel{N3} 04.28 — FCC/motivação}
A Administração aplica sanção a servidor por decisão que apenas afirma ``por interesse público'', sem indicar fatos ou fundamentos. A deficiência principal é de:
\begin{enumerate}[label=\Alph*)]
\item competência legislativa;
\item motivação;
\item personalidade jurídica;
\item publicidade da lei;
\item descentralização.
\end{enumerate}
\end{questao}

\begin{questao}{\nivel{N3} 04.29 — FGV/autotutela}
A Administração descobre ato ilegal próprio e, sem depender de provocação judicial, inicia procedimento para anulá-lo, respeitando contraditório quando necessário. A atuação decorre do princípio da:
\begin{enumerate}[label=\Alph*)]
\item continuidade;
\item autotutela;
\item especialidade;
\item supremacia do Legislativo;
\item indisponibilidade do processo judicial.
\end{enumerate}
\end{questao}

\begin{questao}{\nivel{N3} 04.30 — Cebraspe/Certo ou Errado}
A autotutela administrativa não torna dispensáveis contraditório e ampla defesa quando a invalidação afeta situação jurídica individual consolidada em contexto que exija tais garantias.
\end{questao}

\begin{questao}{\nivel{N3} 04.31 — FCC/serviço público}
A continuidade do serviço público:
\begin{enumerate}[label=\Alph*)]
\item significa impossibilidade absoluta de qualquer interrupção;
\item é princípio relevante, mas admite hipóteses de interrupção previstas no regime jurídico aplicável;
\item elimina o direito de greve em qualquer situação sem mediação constitucional;
\item impede manutenção programada;
\item dispensa planejamento de contingência.
\end{enumerate}
\end{questao}

\begin{questao}{\nivel{N3} 04.32 — FGV/responsabilidade por omissão}
Sobre responsabilidade estatal por omissão, assinale a melhor afirmação.
\begin{enumerate}[label=\Alph*)]
\item Toda omissão gera responsabilidade objetiva automática.
\item O tema exige análise do dever jurídico de agir, do nexo e do regime jurisprudencial aplicável, não cabendo fórmula única para qualquer omissão.
\item O Estado jamais responde por omissão.
\item Só há responsabilidade se houver contrato.
\item O agente responde sozinho.
\end{enumerate}
\end{questao}

\begin{questao}{\nivel{N3} 04.33 — Cebraspe/Certo ou Errado}
A existência de causa exclusiva da vítima pode, conforme o caso, romper o nexo causal e afastar a responsabilidade estatal.
\end{questao}

\begin{questao}{\nivel{N3} 04.34 — FCC/improbidade}
Após a Lei nº 14.230/2021, a análise de improbidade administrativa deve, entre outros aspectos:
\begin{enumerate}[label=\Alph*)]
\item presumir dolo de qualquer ilegalidade;
\item verificar tipificação legal e elemento subjetivo exigido, não equiparando toda ilegalidade a improbidade;
\item aplicar automaticamente a Lei Penal;
\item ignorar prescrição;
\item tratar culpa simples como suficiente em qualquer tipo.
\end{enumerate}
\end{questao}

\begin{questao}{\nivel{N3} 04.35 — FGV/licitações}
Um órgão escolhe fornecedor específico por preferência técnica interna, embora existam vários capazes de competir, e tenta enquadrar a contratação como inexigível. O problema central é:
\begin{enumerate}[label=\Alph*)]
\item inexigibilidade exige inviabilidade de competição demonstrada, não mera preferência;
\item todo fornecedor conhecido gera inexigibilidade;
\item competição é vedada em TI;
\item escolha técnica elimina necessidade de motivação;
\item inexigibilidade e dispensa são sinônimos.
\end{enumerate}
\end{questao}

\begin{questao}{\nivel{N3} 04.36 — Cebraspe/Certo ou Errado}
Na contratação direta, a ausência de licitação não elimina o dever de instrução, motivação, justificativa de preço e demonstração do enquadramento legal.
\end{questao}

\begin{questao}{\nivel{N3} 04.37 — FCC/controle}
O Tribunal de Contas, no exercício do controle externo, integra:
\begin{enumerate}[label=\Alph*)]
\item a Administração direta do Poder Executivo;
\item o sistema constitucional de controle externo exercido pelo Legislativo com seu auxílio;
\item o Poder Judiciário;
\item a função de autotutela do órgão auditado;
\item exclusivamente o controle interno.
\end{enumerate}
\end{questao}

\begin{questao}{\nivel{N3} 04.38 — FGV/supervisão finalística}
Ministério determina diretamente a um presidente de autarquia como decidir processo técnico específico, invocando ``poder hierárquico''. A justificativa é problemática porque:
\begin{enumerate}[label=\Alph*)]
\item autarquia não integra o Estado;
\item não há, em regra, subordinação hierárquica entre pessoas jurídicas distintas; existe vinculação/supervisão nos termos legais;
\item autarquias não estão sujeitas a controle algum;
\item ministérios não podem exercer supervisão;
\item toda autarquia é independente como Poder.
\end{enumerate}
\end{questao}

\begin{questao}{\nivel{N3} 04.39 — Cebraspe/Certo ou Errado]
A discricionariedade pode envolver escolha legítima entre alternativas legais, mas a motivação declarada pode ser controlada quanto à veracidade dos pressupostos fáticos que a sustentam.
\end{questao}

\begin{questao}{\nivel{N3} 04.40 — FCC/Lei 8.112}
A responsabilidade administrativa de servidor:
\begin{enumerate}[label=\Alph*)]
\item exclui sempre a responsabilidade civil e penal;
\item pode coexistir com responsabilidades civil e penal, observadas as relações legais entre as instâncias;
\item somente existe se houver condenação penal;
\item nunca depende de processo disciplinar;
\item é objetiva em qualquer infração.
\end{enumerate}
\end{questao}

\subsection*{N4 — Auditoria / avançado}

\begin{questao}{\nivel{N4} 04.41 — Cebraspe/caso integrado}
Um ato favorável foi praticado há anos com vício, sem indício de má-fé do beneficiário. Julgue o item:

Antes de simplesmente anulá-lo, a Administração deve considerar o regime decadencial e a segurança jurídica aplicáveis, além do contraditório quando pertinente.
\end{questao}

\begin{questao}{\nivel{N4} 04.42 — FGV/controle de mérito}
O Judiciário reconhece que determinado ato discricionário é legal, mas entende que outra escolha seria administrativamente ``mais eficiente''. A providência constitucionalmente mais adequada é:
\begin{enumerate}[label=\Alph*)]
\item substituir a escolha por aquela que o juiz considera melhor;
\item respeitar o espaço legítimo de mérito, salvo vício de juridicidade controlável;
\item revogar o ato;
\item assumir a gestão do órgão;
\item determinar contratação específica sem base legal.
\end{enumerate}
\end{questao}

\begin{questao}{\nivel{N4} 04.43 — FCC/contratação direta}
Auditoria identifica contratação por inexigibilidade com três fornecedores aptos e ausência de demonstração de exclusividade ou inviabilidade de competição. O achado mais adequado é:
\begin{enumerate}[label=\Alph*)]
\item ``Toda inexigibilidade é ilegal'';
\item ``O processo não demonstra a inviabilidade de competição necessária ao enquadramento e requer revisão da motivação, pesquisa de mercado e responsabilização conforme o caso'';
\item ``Basta reduzir o preço'';
\item ``A contratação é válida se o fornecedor for conhecido'';
\item ``Deve-se trocar a modalidade para dispensa sem nova análise''.
\end{enumerate}
\end{questao}

\begin{questao}{\nivel{N4} 04.44 — Cebraspe/Certo ou Errado}
Uma ilegalidade administrativa pode existir sem que, automaticamente, estejam preenchidos todos os requisitos de um ato de improbidade após a reforma da Lei nº 8.429/1992.
\end{questao}

\begin{questao}{\nivel{N4} 04.45 — FGV/poder de polícia}
Órgão aplica restrição a atividade privada sem competência legal e invoca genericamente ``poder de polícia''. O vício principal é:
\begin{enumerate}[label=\Alph*)]
\item poder de polícia afasta legalidade;
\item falta de competência, pois prerrogativa administrativa depende de fundamento jurídico;
\item ausência de personalidade do particular;
\item inexistência de interesse público em qualquer fiscalização;
\item excesso de publicidade.
\end{enumerate}
\end{questao}

\begin{questao}{\nivel{N4} 04.46 — Cebraspe/Certo ou Errado}
A motivação de ato discricionário pode vincular a Administração aos motivos declarados, de modo que a inexistência ou falsidade desses pressupostos seja relevante para o controle de validade.
\end{questao}

\begin{questao}{\nivel{N4} 04.47 — FCC/responsabilidade estatal}
Um terceiro sofre dano decorrente de atuação de empregado de concessionária de serviço público, nessa qualidade. Em tese, o regime constitucional do art. 37, §6º, alcança:
\begin{enumerate}[label=\Alph*)]
\item apenas pessoas jurídicas de direito público;
\item também pessoas jurídicas de direito privado prestadoras de serviços públicos, nos termos constitucionais;
\item exclusivamente o agente pessoa física;
\item somente contratos de obra;
\item nenhuma concessionária.
\end{enumerate}
\end{questao}

\begin{questao}{\nivel{N4} 04.48 — FGV/processo administrativo}
Autoridade indefere recurso administrativo com texto padronizado que não enfrenta nenhum argumento novo e relevante apresentado pelo recorrente. O principal ponto de controle é:
\begin{enumerate}[label=\Alph*)]
\item descentralização;
\item suficiência da motivação e observância do devido processo administrativo;
\item personalidade jurídica do órgão;
\item licitação;
\item responsabilidade objetiva.
\end{enumerate}
\end{questao}

\begin{questao}{\nivel{N4} 04.49 — Cebraspe/Certo ou Errado}
O princípio da eficiência autoriza a Administração a afastar exigência legal sempre que a solução informal for mais rápida e barata.
\end{questao}

\begin{questao}{\nivel{N4} 04.50 — FGV/matriz de achado}
Auditoria identifica renovação contratual por inércia, sem análise de desempenho, vantajosidade ou alternativas, embora o serviço seja crítico. O achado mais adequado é:
\begin{enumerate}[label=\Alph*)]
\item ``Prorrogar contratos é proibido'';
\item ``A ausência de motivação e avaliação de vantajosidade/desempenho fragiliza a decisão de prorrogação e pode perpetuar contratação desvantajosa; recomenda-se processo decisório documentado e análise de alternativas'';
\item ``Deve-se trocar de fornecedor obrigatoriamente'';
\item ``Basta exigir desconto de 1\%'';
\item ``Serviço crítico deve ser sempre internalizado''.
\end{enumerate}
\end{questao}
